% Generated by report_gen.py (ai-ee v1) - do not hand-edit.
\documentclass[11pt,a4paper]{article}
\usepackage[margin=2.2cm]{geometry}
\usepackage{graphicx}
\usepackage{booktabs}
\usepackage{longtable}
\usepackage{pdfpages}
\makeatletter
\@ifundefined{KV@Gin@artifact}{\define@key{Gin}{artifact}[]{}}{}
\makeatother
\usepackage{xcolor}
\usepackage[hidelinks]{hyperref}
\setcounter{tocdepth}{1}
\begin{document}
\sloppy
\begin{center}
{\LARGE\bfseries pd-trigger --- Design Document}\\[6pt]
{\large ai-ee v1 pipeline}\\[2pt]
generated 2026-08-14 09:02:20
\end{center}
\tableofcontents

\section{Board and Run Metadata}
{\small
\begin{longtable}{lp{11cm}}
\toprule
\textbf{Field} & \textbf{Value} \\
\midrule
\endhead
board & pd-trigger \\
workspace & \texttt{boards/pd-trigger} \\
phase & P10 \\
gate status & all 5 recorded gates pass \\
state created & 2026-07-28T05:57:45 \\
state updated & 2026-08-14T01:21:02 \\
\bottomrule
\end{longtable}
}
\section{Overview}
\subsection*{Brief: pd-trigger (S14 run c - novel brief chosen by the user)}
USB-C PD trigger/load board:

\begin{itemize}
\item USB-C receptacle, negotiating USB Power Delivery as a SINK
\item PD sink controller (CH224-class or equivalent)
\item Output voltage selectable among PD profiles (5/9/12/15/20V) via an
onboard selector (DIP switch / jumper / button per architect's choice)
\item Output on a screw terminal (plus an auxiliary 2.54mm header)
\item Status indication: at least power-present; profile indication welcome
\item Input protection appropriate for a bench tool (TVS; fusing if warranted)
\item Power path sized for 5A CONTINUOUS (up to 100W at 20V) - copper, terminal
and connector ratings must reflect this
\item JLCPCB 2-layer, economy PCBA where possible, compact (\textasciitilde{}40x25mm target)
\item Prototype qty 10; JLC Basic parts preferred
\item Bench use, no enclosure, no battery
\end{itemize}
\subsection*{Fabrication spec snapshot}
{\small
\begin{longtable}{ll}
\toprule
\textbf{Key} & \textbf{Value} \\
\midrule
\endhead
layers & 2 \\
width\_mm & 48.0 \\
height\_mm & 30.0 \\
qty & 10 \\
surface\_finish & HASL \\
solder\_mask\_color & green \\
assembly & True \\
\bottomrule
\end{longtable}
}
\subsection*{Stackup}
Chosen stackup: \texttt{JLC2313\_1.6\_2oz} (2-layer, 1.6 mm, \textbf{2 oz} outer copper, HASL)
\section{Requirements}
\subsection*{Requirements: pd-trigger}
Source: \texttt{brief/brief.md} (S14 run c - novel brief). No other attachments. Unstated low-risk items are marked \texttt{ASSUMED:} inline. Section 9 holds every design-changing or safety-relevant unknown - all six items there trace back to the \textgreater{}3A flag (section 8) or a part-selection conflict found while writing this document; none are guessed.

\subsubsection*{1. Function}
A USB-C Power Delivery trigger/breakout board for bench use. A CH224-class (or equivalent) PD sink controller negotiates as a SINK with a connected USB-C PD charger/source, and the user selects one of five PD profiles (5/9/12/15/20V) via an onboard selector (DIP switch, jumper, or button - architect's choice, stated). The negotiated voltage is delivered on a screw terminal (primary output) plus an auxiliary 2.54mm header, sized for 5A continuous (up to 100W at 20V). Status LEDs show at minimum power-present; per-profile indication is optional ("welcome" but not required). No regulation/conversion circuitry beyond the PD negotiation itself - this is a pass-through trigger board, not a DC-DC converter. Bench tool: no enclosure, no battery. Prototype qty 10.

NOTE: delivering 5A continuous through a compact \textasciitilde{}40x25mm 2-layer board is a real copper-sizing constraint (trace/pour width, copper weight, thermal relief) - carried forward as a directive for the architecture stage: verify against a current-capacity calculation, do not assume default 1oz copper is sufficient without checking.

\subsubsection*{2. Interfaces}
\begin{itemize}
\item USB-C receptacle (input): PD SINK only, stated. Negotiates with the
source; no USB data function (no D+/D-, no USB 2.0/3.x data role) implied
or needed - power-only use of the connector. ASSUMED: CC1/CC2 wired per
the chosen PD controller's application circuit (architect's part choice
drives the exact pinout).
\item Output voltage selector: onboard, type left to architect (DIP switch /
jumper / button) - stated explicitly as the architect's choice.
\item Output connector: 2-pole screw terminal (V+/GND), primary output,
stated. Pitch/current rating: see open question 3.
\item Auxiliary output: 2.54mm pin header, stated. Its purpose and current
rating relative to the screw terminal are unstated and matter for part
selection - see open question 4.
\item Status indication: at least 1x power-present LED, stated (mandatory).
Per-profile indication: stated as "welcome," i.e. optional. ASSUMED:
architect adds it only if cheap/simple (e.g. one extra LED or a color
change), omits otherwise - no user decision needed here.
\end{itemize}
\subsubsection*{3. Power}
\begin{itemize}
\item Source: USB-C VBUS, PD-negotiated as SINK only. No battery, no mains, no
other input (stated: "no battery").
\item Negotiated profiles: 5V / 9V / 12V / 15V / 20V, user-selected (stated) -
the profile set CH224-class trigger controllers natively support
(matches the brief exactly).
\item Design current: 5A continuous per the brief, stated as sizing "the power
path" generally (copper, terminal, connector). Whether that means a
uniform 5A at EVERY profile (25W at 5V up to 100W at 20V) or a
per-profile-varying cap is the single most consequential unknown in this
document - see open question 1.
\item 100W at 20V (5A x 20V) is within normal USB PD range but only reachable
if the source advertises a 20V/5A PDO AND the cable is electronically
marked (e-marked) for 5A - both outside this board's control. Noted here
as a fact, not a question: the board must be built to handle 5A when
offered, not guaranteed to receive it from every source/cable pairing.
\item USB-C receptacle current rating: the connector itself is spec'd for up
to 5A (with an e-marked cable) by the USB-C standard - no separate
rating decision needed for the input connector.
\item Rail budget for the controller/logic/LEDs (GUESS): low tens of mA,
negligible next to the 5A main path. No separate regulator is implied by
the brief (pass-through trigger board, not a DC-DC converter). ASSUMED:
the PD controller supplies its own housekeeping rail internally (typical
for CH224-class parts) or the selector/LEDs draw directly off VOUT;
architect confirms against the chosen part.
\item Fusing / input protection: TVS is stated (add it). Fusing is explicitly
conditional in the brief ("if warranted") - for a \textgreater{}3A path this is
safety-relevant and not guessable - see open question 2.
\end{itemize}
\subsubsection*{4. Environment}
Not stated beyond "bench use, no enclosure" (stated). ASSUMED: indoor lab/bench ambient (roughly 0-40C), no ingress protection, no vibration/shock requirement. Low risk - qty-10 prototype, easy to revisit.

\subsubsection*{5. Size \& mounting}
\begin{itemize}
\item Outline: brief says "compact (\textasciitilde{}40x25mm target)" - a target, not stated
as a hard max. Fitting a USB-C receptacle, a 5A-rated screw terminal, PD
controller, selector, and LEDs is workable at this size but tight, and
the tolerance if the layout runs slightly over is unstated - see open
question 5.
\item Mounting holes: not stated. ASSUMED: none (bench use, no enclosure).
\item Height limit: not stated. ASSUMED: none (open board, no enclosure); the
screw terminal and USB-C receptacle heights are the practical limits.
\end{itemize}
\subsubsection*{6. Quantity \& budget}
\begin{itemize}
\item Build quantity: prototype qty 10 (stated).
\item Target unit cost: not stated. ASSUMED: no hard cap; minimized via the
JLC Basic parts preference + economy PCBA (stated strategy).
\end{itemize}
\subsubsection*{7. Assembly}
\begin{itemize}
\item JLCPCB, 2-layer, economy PCBA "where possible" (stated). ASSUMED
interpretation: JLC assembles everything it can from its catalog; any
part that isn't JLC-catalog-assemblable (e.g. the screw terminal,
possibly the USB-C receptacle depending on footprint/stock) is
hand-soldered post-PCBA; architect decides per actual part availability.
\item Parts: JLC Basic parts preferred (stated), not exclusive - Extended or
hand-solder parts allowed where Basic doesn't cover the need (e.g. a
5A-rated screw terminal, the PD controller itself).
\item Sidedness: not stated. ASSUMED: single-sided (top) SMT preferred for the
economy tier and small part count; architect may go double-sided if
placement on \textasciitilde{}40x25mm requires it.
\end{itemize}
\subsubsection*{8. Compliance/safety flags}
\begin{itemize}
\item Mains voltage: no - USB-C PD input only, max 20V.
\item Batteries: no (stated).
\item Motors: no - output is a generic power tap, load-agnostic.
\item \textgreater{}30V: no - max 20V from PD negotiation.
\item High current (\textgreater{}3A): YES - power path is stated to be sized for 5A
continuous, up to 100W at 20V. This is the primary safety driver for
this board and the source of open questions 1-4 below: nothing about the
current path is being guessed.
\item RF transmit: no - no radio/wireless function.
\end{itemize}
\subsubsection*{9. Open questions}
All six items below stem from the \textgreater{}3A safety flag (section 8) or from a part-selection conflict found while writing this document. The pipeline will not proceed on assumptions here.

1. Max current commitment: does 5A continuous apply UNIFORMLY to every selectable profile (copper, terminal, and connector all handle 5A even at 5V/9V/12V/15V, not just at 20V), or does the design current vary by profile (lower at lower voltages)? DEFAULT: uniform 5A at every profile - simplest, safest, matches the brief's unconditional wording ("power path sized for 5A continuous"), and gives headroom since a typical PD source will offer less than 5A at the lower voltages anyway.

2. Fusing: add a fuse/current-limiting element on the input, and if so what type? The brief leaves this conditional ("fusing if warranted") but a 5A/100W path is exactly the case where it usually is. Options: (a) none - rely on the PD source's own protections plus the TVS only; (b) resettable PTC/polyfuse (survives a downstream/output short without replacement); (c) one-time fast-blow fuse. DEFAULT: (b) resettable PTC, sized to hold at the committed current (question 1) with headroom and trip well above it - cheap, JLC-Basic- friendly, and self-resetting after a bench-tool short-circuit mistake.

3. Screw terminal sizing: what wire gauge should the output terminal accept, and what minimum current rating (component nameplate rating, not just the 5A operating point) should it carry? These two drive the same part choice (pitch/size), so they're answered together. Options: (a) 20AWG / rated \textgreater{}=5A (smallest footprint, least margin); (b) 18AWG / rated \textgreater{}=8A (comfortable margin); (c) 16AWG / rated \textgreater{}=10A (most margin, largest footprint - pressures the size target in question 5). DEFAULT: (b) 18AWG, terminal rated \textgreater{}=8A - standard mid-size terminal block class, reasonable margin above 5A without oversizing the board.

4. Auxiliary 2.54mm header purpose: standard 2.54mm pin headers are typically only rated \textasciitilde{}1-3A PER PIN - well under the 5A design current - unless multiple pins are paralleled per net. The brief lists this header alongside the screw terminal without saying which duty it carries. Options: (a) full-power tap, same 5A as the screw terminal - requires paralleling 3+ pins per net (VOUT, GND), extra part count/footprint; (b) low-current/sense tap only (e.g. probing VOUT, driving small accessory logic) - NOT rated for the full 5A; (c) full-power tap via a beefier high-current 2.54mm-pitch connector instead of a plain header. DEFAULT: (b) - auxiliary header is a low-current tap only; the screw terminal is the sole full-current output. Simplest and smallest, and avoids silently under-rating a 5A path onto what looks like a generic 0.1in header.

5. Size target tolerance: "\textasciitilde{}40x25mm target" - is this a hard ceiling (redesign if exceeded) or a soft target (get as close as practical, given a 5A-rated screw terminal, USB-C receptacle, and selector all need physical room)? DEFAULT: soft target - architect minimizes toward \textasciitilde{}40x25mm but may run modestly over (roughly up to 15-20\%) if needed to fit the current-rated connector/terminal safely rather than undersizing them to hit the number.

6. PD negotiation fallback behavior: CH224-class controllers commonly fall back to 5V (rather than holding or erroring) when the source can't deliver the user-selected profile (e.g. a weak/cheaper charger, or a non-PD source). Is a silent 5V fallback on the output terminal acceptable to whatever load the user connects, or must the board react differently? Options: (a) yes, silent fallback is fine - board just outputs whatever gets negotiated; (b) no - actively block/disconnect the output (e.g. via a load switch) unless the SELECTED profile is confirmed, rather than ever presenting an unrequested voltage; (c) fallback is fine but must be visibly indicated (status LED distinguishes "selected profile achieved" from "fallback/mismatch"). DEFAULT: (c) - keeps the design simple (no added load switch) while making a wrong-voltage condition visible, which matters for a bench tool where the user picks a profile expecting the load to see that voltage.

\subsubsection*{Answers (P0, decided by orchestrator as user's delegate per directive)}
1. Uniform 5A at every profile. 2. Resettable PTC fuse. 3. Screw terminal 18AWG-capable, \textgreater{}=8A rated. 4. Aux header LOW-CURRENT tap only (silk-marked). 5. Size soft target (up to \textasciitilde{}20\% over allowed for current-rated parts). 6. 5V fallback ACCEPTABLE but must be VISIBLY indicated (profile LED scheme).

\section{Architecture}
\subsection*{Decisions of record (state.json)}
{\small
\begin{longtable}{lp{5.2cm}p{7.2cm}}
\toprule
\textbf{Phase} & \textbf{What} & \textbf{Why} \\
\midrule
\endhead
P0 & P0 safety answers as delegate: uniform 5A, PTC fuse, 8A screw terminal, aux header low-current only, soft size, indicated 5V fallback & All defaults engineering-sound; \textgreater{}3A flag addressed explicitly per analyst framing \\
P1 & P0 answer 2 REVISED: NO main-path series fuse (PD source OCP + mandatory TVS are the protection; no SMD PPTC holds 5A@40C and trip window unreachable under source foldback); aux tap gets 1A 1206 PPTC + silk AUX 1A MAX & Two independent research fragments converged (interface + power); THT 24x25mm disc would break size and add hand-solder for negative value \\
P1 & 2oz outer copper: stackups.yaml gained JLC2313\_1.6\_2oz (S14 reference fix); 5A width 1.75mm vs unroutable 3.5mm at 1oz & Power research IPC-2152 math via the pipeline's own model \\
P2 & Flat root sheet; VBUS one net end-to-end; GND out of power{[}{]} (structural pour + 10-via rule); DIP straps fail-safe-to-5V; zener+dual-NPN window LED scheme for fallback indication & Architect P2 with all 10 settle-items reasoned \\
P4 & Retroactive approval: P4 edited architecture/constraints.json (R2-\textgreater{}R2A/R2B, C1-\textgreater{}C1A/C1B in separation/groups) - place\_anneal silently drops absent refs, stale refdes = invisible constraint loss (finding 31) & Correct and mandated by sheets.md tracking rule \\
P7 & VBUS fan-in = F.Cu pour (measured: no 1.75mm track can touch J1 VBUS pads, max 1.465mm); trunk 3.0mm track; 0 VBUS vias; C1B bulk tap via 2 pour channels 1.11+1.25mm (disclosed - no checker measures viasless pour channels); netclass split in .kicad\_pro only (rules\_gen one-class defect) & Router-measured geometry falsified the placement premise; zones are not width-rule-bound; DRU authority untouched \\
P8 & Outline 48x30 vs D9 ceiling 48x28: DEVIATION ACCEPTED (2mm on soft axis; P0 literal +20\% of 25 = 30; no outline-edit op exists at v1) & Reviewer-caught undisclosed drift, now disclosed; re-initing P5 for 2mm not warranted \\
P8 & C1B pour channels 1.11+1.25mm = STANDING WAIVER (dead-end bulk stub, 0.3A transient vs 3.49A capacity) & Reviewer accepted with numbers \\
\bottomrule
\end{longtable}
}
\subsection*{Sheet plan}
\subsection*{pd-trigger - sheet plan}
{*}Amended A1 after the P3 datasheet extracts: \texttt{+3V3} and U2 are gone, \texttt{/VIND} and \texttt{/VDD} are new, PG is unconnected. See \texttt{decisions.md} section A1.{*}

\subsubsection*{One flat root sheet - \texttt{pd-trigger.kicad\_sch}}
\begin{quote}\small\ttfamily\raggedright
\textbar{} Sheet \textbar{} File \textbar{} Blocks \textbar{} Interface nets (sheet pins) \textbar{} pwr\_base \textbar{}\\
\textbar{}---\textbar{}---\textbar{}---\textbar{}---\textbar{}---\textbar{}\\
\textbar{} root `pd-trigger` \textbar{} `kicad/pd-trigger.kicad\_sch` \textbar{} B1-B6, all of them \textbar{} {*}{*}none{*}{*} - there are no child sheets \textbar{} {*}{*}1{*}{*} \textbar{}
\end{quote}
Justification (the plan allows a flat sheet if argued):

\begin{itemize}
\item 28 components in one functional chain. The usb-buck precedent split at 28
components across three genuinely separable domains; this board has one.
\item \textbf{Every net that crosses a block boundary is a rail or a two-part stub}
(\texttt{VBUS}, \texttt{GND}, \texttt{/VIND}, \texttt{/VDD}, \texttt{/VAUX}). Power symbols are global and need no
sheet pin, so hierarchy would add sheet pins carrying nothing while introducing
\texttt{/{[}sheet{]}/NAME} prefixes into net names that \texttt{constraints.json} would track.
\item The only multi-block signal is \texttt{/HV\_OK}, inside the indication block's own
fan-out.
\item P4 parallelism is not worth buying here: one generator script, one ERC run,
one netlist audit.
\end{itemize}
Consequence for P4: \textbf{all local labels become `/NAME`} (root-local), which is what \texttt{constraints.json} assumes for \texttt{/VAUX}, \texttt{/VIND} and \texttt{/VDD}.

\subsubsection*{1. Canonical net names}
These are contractual. \texttt{check\_current} raises \texttt{CheckError} (exit 2) on a \texttt{power} net that is absent from the board, so \texttt{VBUS}, \texttt{/VAUX}, \texttt{/VIND} and \texttt{/VDD} must appear verbatim in the final netlist.

\paragraph*{Global rails - power symbols, bare names}
\begin{quote}\small\ttfamily\raggedright
\textbar{} Net \textbar{} Symbol \textbar{} Driven by \textbar{} PWR\_FLAG? \textbar{}\\
\textbar{}---\textbar{}---\textbar{}---\textbar{}---\textbar{}\\
\textbar{} `VBUS` \textbar{} `power:VBUS` \textbar{} J1 VBUS pins (passive) \textbar{} {*}{*}yes{*}{*} \textbar{}\\
\textbar{} `GND` \textbar{} `power:GND` \textbar{} passive everywhere \textbar{} {*}{*}yes{*}{*} \textbar{}
\end{quote}
\textbf{`+3V3` NO LONGER EXISTS.} The 3.3 V LDO was removed at amendment A1: CH224K's VDD is an internally shunt-regulated node (3.24/3.30/3.36 V, 0-30 mA sink) that an external regulator would fight. There is no logic rail on this board.

\textbf{`VOUT` does not exist either.} The board is a pass-through: the receptacle VBUS pins, the TVS, the bulk caps, every tap and the screw terminal are {*}{*}one copper object{*}{*}, named \texttt{VBUS} end to end. The interface research's separate \texttt{VOUT} entry (which assumed a series PPTC that no longer exists) is deduplicated INTO \texttt{VBUS}; do not emit two power entries for one net.

\paragraph*{Root-local nets - \texttt{/NAME}}
\begin{quote}\small\ttfamily\raggedright
\textbar{} Net \textbar{} Members \textbar{} Note \textbar{}\\
\textbar{}---\textbar{}---\textbar{}---\textbar{}\\
\textbar{} `/VAUX` \textbar{} F1 pin 2, J3 pin 1 \textbar{} own net so `check\_current` sizes the aux stub at 1 A, not 5 A \textbar{}\\
\textbar{} `/VIND` \textbar{} R14 pin 2, R1, R2, R6, R8, R10, R12, R13 \textbar{} housekeeping stub behind the 0 ohm link; `pdn: false`, width-only \textbar{}\\
\textbar{} `/VDD` \textbar{} R2 pin 2, U1 pin 1 (VDD), C5, R3/R4/R5 \textbar{} the shunt node. Needs a {*}{*}PWR\_FLAG{*}{*} - U1 pin 1 is a power-input pin fed through a passive \textbar{}\\
\textbar{} `/CFG1` `/CFG2` `/CFG3` \textbar{} U1 pins 9/2/3, R3/R4/R5, SW1 \textbar{} 100 k pull-ups to {*}{*}`/VDD`{*}{*}, switch shorts to `GND` \textbar{}\\
\textbar{} `/VSENSE` \textbar{} R1, U1 pin 8 \textbar{} 10 k series is mandatory (13.5 V abs max pin on a 21 V rail) \textbar{}\\
\textbar{} `/CC1` `/CC2` \textbar{} J1 A5 -\textgreater{} U1 pin 7, J1 B5 -\textgreater{} U1 pin 6 \textbar{} straight through, {*}{*}no crossover, no series R, no external Rd{*}{*} \textbar{}\\
\textbar{} `/BC12\_DIS` \textbar{} U1 pin 4 (DP) + pin 5 (DM) tied together \textbar{} PD-only operation - {*}{*}one net, deliberately shorted{*}{*} (see V12) \textbar{}\\
\textbar{} `/ZBIAS` \textbar{} D2 anode, R6 pin 1 \textbar{} window detector bias chain from `/VIND` \textbar{}\\
\textbar{} `/HV\_B` \textbar{} R6 pin 2, Q1A base, R7 \textbar{} R7 4k7 to GND kills zener knee leakage \textbar{}\\
\textbar{} `/HV\_OK` \textbar{} Q1A collector, R8, D6 cathode, R9 \textbar{} high above \textasciitilde{}6.7 V on the bus = profile achieved \textbar{}\\
\textbar{} `/FB\_B` \textbar{} R9 pin 2, Q1B base \textbar{} inverter drive \textbar{}\\
\textbar{} `/FB\_K` \textbar{} Q1B collector, D5 cathode \textbar{} red "5V ONLY" leg \textbar{}
\end{quote}
LED anode nets (R10/R12/R13 to D3/D5/D6) are two-part local nets; leave them auto-named. \textbf{U1 pin 10 (PG) is a no-connect} - see s3.

\textbf{`/BC12\_DIS` naming is load-bearing.} \texttt{constraints.json} sets an explicit \texttt{"diff\_pairs": {[}{]}}, which disables \texttt{check\_diffpair} outright - but the net name also deliberately avoids any \texttt{DP}/\texttt{DM}/\texttt{\_P}/\texttt{\_N}/\texttt{D+}/\texttt{D-} token so that even auto-discovery could not construct a "pair" out of a node that is a single deliberate short.

\subsubsection*{2. Refdes allocation (one namespace, one sheet)}
\begin{quote}\small\ttfamily\raggedright
\textbar{} Block \textbar{} Refs \textbar{} Parts \textbar{}\\
\textbar{}---\textbar{}---\textbar{}---\textbar{}\\
\textbar{} B1 input + protection \textbar{} J1, D1, C1, C2 \textbar{} receptacle, TVS, 22 uF bulk, 100 nF \textbar{}\\
\textbar{} B2 controller + supply \textbar{} U1, R1, {*}{*}R2{*}{*}, C5 \textbar{} CH224K, 10 k sense, {*}{*}1 k / 2512 dropper{*}{*}, 1 uF on /VDD \textbar{}\\
\textbar{} B3 profile selector \textbar{} SW1, R3, R4, R5 \textbar{} DIP-3, 3 x {*}{*}100 k{*}{*} pull-up to /VDD \textbar{}\\
\textbar{} B4 housekeeping stub \textbar{} {*}{*}R14{*}{*} \textbar{} 0 ohm 0603 link, VBUS -\textgreater{} /VIND \textbar{}\\
\textbar{} B5 indication \textbar{} D2, Q1, R6, R7, R8, R9, D3, D5, D6, R10, R12, R13 \textbar{} zener, dual NPN, window network, 3 LEDs + legs \textbar{}\\
\textbar{} B6 output + aux \textbar{} J2, F1, J3 \textbar{} screw terminal, 1 A PPTC, aux header \textbar{}\\
\textbar{} power symbols \textbar{} `\#PWR01`+ / `\#FLG01`+ \textbar{} pwr\_base = 1 \textbar{}
\end{quote}
\textbf{Refdes gaps are intentional: there is no U2, C3, C4, D4 or R11.} Those were the LDO, its two capacitors, the PG LED and the PG LED's resistor, all removed at amendment A1. Surviving parts keep their original designators so that P3's in-flight sourcing does not have to be re-keyed - do \textbf{not} renumber to close the gaps.

Assignments are reflected in \texttt{constraints.json} (\texttt{thermal} names J1; \texttt{placement.edges} names J1, J2, J3, SW1; \texttt{groups} name Q1 and D3; \texttt{separation} names R2 and U1). {*}{*}\texttt{check\_thermal} hard-errors (exit 2) on a refdes with no pads{*}{*}, so if P4 renumbers, \texttt{constraints.json} must be renumbered with it.

\subsubsection*{3. Wiring facts P4 must not re-derive from memory}
All from \texttt{parts/C970725.json} (CH224 manual V2.1) unless noted:

1. \textbf{Pin 0 is the exposed baseplate AND the ground terminal}, and \textbf{KiCad/EasyEDA footprints usually number that pad 11}. Run \texttt{schlib.py --pins} against the actual symbol and check the footprint's pad numbering before wiring; a mis-numbered thermal pad leaves GND unconnected on a part whose only ground is that pad. 2. \textbf{No external 5.1 k Rd on CC1/CC2.} The CH224K reference schematic wires CC straight to the connector; the CH224D and CH221K schematics in the same manual \emph{do} show 5.1 k. Rd is integrated on this part - fitting external pull-downs would put \textasciitilde{}2.55 k on CC and shift the source's Rd/Ra detection. 3. \textbf{CC1 = A5 -\textgreater{} pin 7, CC2 = B5 -\textgreater{} pin 6, no crossover, no series elements} (a series R adds directly to Rd, whose budget is only +/-20 \%). 4. \textbf{No caps on CC} without knowing U1's own pin capacitance (the PD spec's CC receiver window is 200-600 pF total, and 200 pF is a \emph{minimum}). 5. \textbf{R2 (1 k) into pin 1 and R1 (10 k) into pin 8 are mandatory.} Pin 1 is a 3.0-3.6 V shunt node; pin 8 is a 13.5 V-max detect input on a 21 V rail. Neither may be shorted to the bus, and neither may be "optimised" away. 6. \textbf{CFG straps: 100 k to `/VDD`, switch to GND.} CFG2/CFG3 have {*}{*}no internal pull-ups{*}{*} and are absolute-max VDD + 0.5 V - they must never see VBUS. 10 k would eat 58 \% of the VDD budget at the 5 V profile. 7. \textbf{PG (pin 10) is left unconnected}, matching the datasheet's reference schematic. It needs an explicit \textbf{no-connect flag} with the justification comment: {*}"open-drain, no absolute-maximum rating published for CH224K - may not be pulled to VBUS/VIND; a VDD-referenced pull-up would spend the 1.7 mA VDD budget. Indication is handled by the D5/D6 window (blocks.md B5)."{*} 8. \textbf{DP (pin 4) and DM (pin 5) shorted together at the chip, off the connector} (\texttt{/BC12\_DIS}); the receptacle's A6/A7/B6/B7 stay unconnected. See \textbf{V12} - the extract reads the reference schematic as wiring DP/DM to the connector, while the interface research cites a PD-only-mode instruction to short them. Resolve before wiring; do not leave them floating either way. 9. \textbf{J1 shell tied directly to GND} (bench tool, no chassis - no 1 M / 4.7 nF hybrid tie). 10. \textbf{All four VBUS pins and all four GND pins of J1 in the netlist} - the 5 A rating is collective across them, and \texttt{netlist\_audit} is the check that they were not partially wired. 11. Every unconnected receptacle pin (SBU1/2, DP/DM, any SuperSpeed pads) gets an explicit no-connect flag with a one-line justification.

\subsubsection*{4. Decoupling metadata P8 depends on}
\texttt{check\_pdn} reads \texttt{decoupling.json} associations and errors on a {*}{*}declared power rail with no associated cap at all{*}{*}. Two declared nets carry \texttt{pdn} true:

\begin{quote}\small\ttfamily\raggedright
\textbar{} Net \textbar{} Cap that must be associated \textbar{} Value \textbar{}\\
\textbar{}---\textbar{}---\textbar{}---\textbar{}\\
\textbar{} `VBUS` \textbar{} C1 (bulk) and C2 \textbar{} 22 uF, 100 nF \textbar{}\\
\textbar{} `/VDD` \textbar{} {*}{*}C5{*}{*}, associated to {*}{*}U1 pin 1{*}{*} \textbar{} 1 uF \textbar{}
\end{quote}
C5 is the datasheet's only specified capacitor for this part ("VDD: external 1 uF capacitor to GND, series resistor to VBUS"), and declaring \texttt{/VDD} in \texttt{power{[}{]}} is what makes \texttt{check\_pdn} enforce it. Pass \texttt{rail\_net="/VDD"} to \texttt{place\_ic\_with\_decoupling} so the metadata records the final net name rather than a wiring label.

\texttt{/VAUX} and \texttt{/VIND} both carry \texttt{"pdn": false} and are skipped by the inventory: nothing decouples a fused stub to a header, or a resistor-fed indicator stub, by design.

Full architecture narrative (not inlined; see the artifact index): \texttt{architecture/blocks.md}, \texttt{architecture/decisions.md}, \texttt{architecture/power\_tree.md}, \texttt{architecture/stackup.md}.
\section{Schematic}
Gate \texttt{erc} (phase P4): \textbf{pass} after 1 attempt(s); last 2026-07-28T08:49:33, failing 0 of 0.
\subsection*{Review waivers}
\subsection*{Schematic review waivers (pd-trigger)}
W1 dp-dm-short-at-chip: architecture's PD-only short (/BC12\_DIS) ACCEPTED by the reviewer as hardware-SAFER than the datasheet's connector wiring (removes a 21V-onto-3.8V-abs-max fault path). Cost: no QC/BC legacy support - the board is "PD sources only"; goes on B.SilkS + order docs. W2 vdd-worst-corner: true worst case 0.979mA at 4.4V low-line all-straps- closed vs unpublished CH224K IDD - bounded by WCH's own reference circuit; V1 bring-up step: measure VDD at the 5V profile; contingency = both 510R -\textgreater{} smaller (documented floor 600R total). Bring-up addenda: buzz out Q1 E/B/C once (clone PDF JS-gated; Nexperia- standard pinning verified as the compatibility claim).

The full schematic PDF follows.
\includepdf[pages=-]{reports/schematic.pdf}
\section{Layout}
Gate \texttt{place} (phase P6): \textbf{pass} after 1 attempt(s); last 2026-07-28T10:14:30, failing 0 of 0.

Gate \texttt{drc\_routed} (phase P8): \textbf{pass} after 4 attempt(s); last 2026-07-28T11:55:26, failing 0 of 0.
\begin{center}
\includegraphics[width=0.72\textwidth]{reports/render_final/top.png}\\
{\small\texttt{reports/render\_final/top.png}}
\end{center}
\begin{center}
\includegraphics[width=0.72\textwidth]{reports/render_final/bottom.png}\\
{\small\texttt{reports/render\_final/bottom.png}}
\end{center}
\begin{center}
\includegraphics[width=0.95\textwidth]{reports/layers/pd-trigger_layers.png}\\
{\small\texttt{reports/layers/pd-trigger\_layers.png}}
\end{center}
\section{Verification}
Gate \texttt{verify} (phase P8): \textbf{pass} after 3 attempt(s); last 2026-07-28T11:55:27, failing 0 of 0.
\subsection*{Verification checks}
{\small
\begin{longtable}{llcp{6cm}}
\toprule
\textbf{Check} & \textbf{Status} & \textbf{Findings} & \textbf{Note} \\
\midrule
\endhead
\texttt{check\_return\_path} & skipped & 0 & skipped - no constraints \\
\texttt{check\_current} & skipped & 0 & skipped - no constraints \\
\texttt{check\_decoupling} & skipped & 0 & skipped - no decoupling \\
\texttt{check\_diffpair} & pass & 0 &  \\
\texttt{check\_creepage} & skipped & 0 & skipped - no constraints \\
\texttt{check\_thermal} & skipped & 0 & skipped - no constraints \\
\texttt{check\_silk} & pass & 0 &  \\
\texttt{check\_pdn} & skipped & 0 & skipped - no constraints/decoupling \\
\bottomrule
\end{longtable}
}
\subsection*{Design review of record}
\subsection*{Board review (P8 adversarial) - pd-trigger}
Scope: hunt for what the green gates cannot see. Board \texttt{kicad/pd-trigger.kicad\_pcb} (routed, drc\_routed 0/0 err/warn, verify 8/8 checks 0 findings). Methods: renders (\texttt{reports/render\_routed/}), raw-copper rasters from the board file (validated against plane\_repair's own area figure: 1322 vs 1307 mm2), pad/segment/via extraction, and the pipeline's own IPC-2152 model (\texttt{check\_current.required\_width\_mm}) for every neck. Settled items (fuse policy, pour legality, netclass split, LED ceiling, DIP choice) not re-opened.

\subsubsection*{Verdict summary}
The 5 A FORWARD path is genuinely good: solid pour fan-in at J1, one 3.0 mm F.Cu trunk (model: 4 C rise at 5 A), zero VBUS vias, taps compliant. The board's two real problems are (1) the 5 A RETURN transition at J1 - the architecture's own "\textgreater{}= 10 vias at J1's GND pads" substitute for the undeclared-GND blind spot was never implemented; what is there is two 0.2 mm trace necks and a single signal via, and (2) the entire functional silkscreen package promised by the architecture (profile table, AUX 1A MAX, PD-only, LED legends, polarity) does not exist - the board file contains zero gr\_text; the only text anywhere is refdes. check\_silk gates legibility of existing silk only; nothing machine-checks silk content, and constraints.json carries no silk key.

\subsubsection*{Findings (worst first)}
\paragraph*{E1 - 5 A return chokes at J1's GND contact pads (error, checker-blind by design D7)}
The 4 GND cable contacts land on two merged SMD pads (F.Cu only): A1-B12 at (17.63,39.31) and B1-A12 at (17.63,45.71), nominally 2.5 A each at the 5 A rating. Their ONLY copper into the B.Cu return pour:

\begin{itemize}
\item A1-B12 (top): a 0.2 mm wide, 2.7 mm long F.Cu chain (3 segments) into shell pad 1's
THT barrel, plus one 0.6/0.3 via whose annulus grazes the pad corner by \textasciitilde{}0.05 mm
(etch-tolerance-level). 0.2 mm at 2 oz = 0.80 A capacity at 10 C rise (pipeline
model); at 2.5 A the same model gives \textasciitilde{}133 C rise (short-neck heatsinking will
moderate this, but it is a \textgreater{}3x overcurrent on the model's own terms). \textasciitilde{}13.6 squares
= 3.3 mohm.
\item B1-A12 (bottom): a 0.2 mm x 1.3 mm trace into ONE 0.6/0.3 via. The board's own via
budget (constraints via\_amps 0.5) says 2.5 A needs 5 vias; there is 1. Shell pad 4
sits 0.32 mm away but is NOT connected on F.Cu - its copper does not help.
\end{itemize}
Architecture D7/power\_tree structural requirement "\textgreater{}= 10 vias at J1's GND pads" (the stated substitute for keeping GND out of power{[}{]} so check\_current stays quiet): not delivered. DRC cannot see it (0.2 mm is netclass-legal), check\_current cannot see it (GND undeclared), check\_return\_path no-ops (no high\_speed key). Failure mode: sustained 5 A at any profile heats the two necks far past the 10 C design intent; long-term trace/joint fatigue, intermittent ground, and full return current cascading onto the surviving 0.2 mm chain. Fix is cheap at P9: an F.Cu GND spandrel bonding each contact pad to its shell pad (the VBUS pour top edge is at y 39.7 - the strip above it, and the strip x \textless{} 17.2 plus the region south of y 46, are free) plus 4-5 vias per side.

\paragraph*{E2 - Safety/waiver silk markings absent (error, promised artifact)}
Board-wide gr\_text count is ZERO; B.SilkS carries nothing at all (both B.SilkS references in the file are the stackup layer table). Missing, each one a written promise:

\begin{itemize}
\item \texttt{AUX 1A MAX} + \texttt{V+}/\texttt{GND} at J3 (P0 answer 4 "silk-marked", V10). F1 makes overload
safe-by-hardware, but polarity of the 2-pin aux header is unmarked.
\item \texttt{PD SOURCES ONLY} (ERC waiver W1's standing condition: "goes on B.SilkS + order
docs"). The waiver's silk half is currently unfulfilled.
\item J2 screw-terminal polarity. Trunk lands on J2 pad 1 = the LOWER screw (board-bottom
side); nothing anywhere tells the bench user which screw is V+. Reverse-wiring a
load on a 100 W output is the single easiest field mistake this board can cause.
\end{itemize}
\paragraph*{E3 - User-interface silk absent (error, promised artifact)}
\begin{itemize}
\item The five-row profile table on B.SilkS - D3 calls silk "the user interface" and
specifies the table verbatim; the bottom is a clean pour with room. Without it the
ON/OFF-to-voltage map exists only off-board; a user WILL set 9 V expecting 12 V by
trial and error. (Fail-safe floor: any wrong guess \textgreater{}= the 5 V fallback, and OFF-fail
= 5 V - the mistake is bounded but real.)
\item LED legends \texttt{PWR} / \texttt{5V ONLY} / \texttt{PROFILE OK} (D4/V10): the render shows refdes
D3/D5/D6 only. The entire P0-answer-6 "visible fallback indication" scheme is mute
without the words - a lit red LED labeled "D5" indicates nothing.
\item Cable disclaimer "20V @ 5A needs a 5A e-marked cable and a 100W source" (V10 lists
it as silk; power\_tree s7 allows documentation - at minimum it must reach the order
docs at P9/P10, which do not exist yet; fab/ is empty).
\item PRESENT and correct: SW1's own footprint prints \texttt{ON} + \texttt{1 2 3} (the compact F.SilkS
marker D3 asked for) - confirmed on the top render.
\end{itemize}
\paragraph*{W1 - J1 mouth recessed \textasciitilde{}0.4-0.5 mm behind the board edge (warning, flagged at P6)}
Courtyard/silk front edge x = 10.18 vs board edge x = 9.8. USB-C receptacles are normally flush-to-overhanging; a recessed mouth lets thick/molded plug overmolds butt the board edge and hold the plug fractionally short of latch - intermittent CC, no PD, "works with some cables". On a bare bench board with arbitrary cables this is a real annoyance risk. Fix: nudge J1 -x 0.5-1.0 mm at P9 (pour edge follows; nothing else in that corridor), or accept and note in order docs.

\paragraph*{W2 - Outline 48 x 30 mm exceeds the architecture's hard ceiling (warning, drift)}
Edge.Cuts = (9.8,27.51)-(57.8,57.51) = 48 x 30 mm. D9: nominal 45 x 25, "hard ceiling 48 x 28 mm"; 30 mm violates D9's own ceiling by 2 mm even though it equals P0 answer 5's literal +20 \% (40x25 -\textgreater{} 48x30, +44 \% area vs the brief's target). No digest discloses the growth past D9. Consequence is cost/footprint only, but it must be disclosed to the user at the checkpoint rather than discovered.

\paragraph*{A1 - J2 GND pad ties to the return pour through 4 thermal spokes (advisory)}
J2 pad 2 (THT, 2.4 mm pad / 1.6 mm drill) connects via 4 x 0.5 mm diagonal spokes = 2.0 mm total at 2 oz -\textgreater{} model 8 C rise at 5 A. Legal and functional, but these are the most-loaded copper necks on the whole return after E1, at a hand-soldered joint. Recommend a solid-connect override (or a short fat GND strap) at P9; zero cost. (Verified all 4 spokes present; the VBUS pad 1 island below it is barrel-tied, normal.)

\paragraph*{A2 - F1's VBUS feed branch is 18.3 mm long (advisory, V8 letter-violation)}
V8/D8 wanted all five VBUS taps hugging the run with 2-3 mm stubs. D1, C1A, C1B, C2, R14 comply; F1's feed runs 8.4 + 6.5 + 3.4 mm at 1.75 mm width to the top-left corner. Width-compliant (V8's actual purpose), carries \textless{}= 1 A in service; only cost is \textasciitilde{}30 mm2 of extra always-live 20 V copper. No action needed; recorded as drift.

\paragraph*{A3 - Floating 0.65 mm2 GND sliver at (20.0, 28.5) (advisory, cleanup)}
plane\_repair's own dead-island fact (0 anchors), between J3/SW1 clearances. Cosmetic / minor DFM nit; remove at P9 cleanup or accept.

\paragraph*{A4 - /VAUX at PPTC max-non-trip current (advisory)}
/VAUX routes at 0.3 mm (required 0.25 at the declared 1.0 A -\textgreater{} 7 C). F1's worst sustained NON-trip current is \textasciitilde{}1.8 A -\textgreater{} model 27 C rise on those runs until/unless the PPTC trips. Bounded, no damage threshold approached (70-80 C copper on a 20 V node); acceptable. Recorded so nobody "fixes" it blind later.

\subsubsection*{Disclosed-item ruling (hunt item 1): C1B pour channels - ACCEPTED}
P7's disclosure: C1B's tap reaches the VBUS pour bottom lobe through two pour channels 1.11 + 1.25 mm wide (between D1's GND-pad keepouts and the stitch-via keepouts) - checker-blind because zones are exempt from track width rules. Judgment with numbers:

\begin{itemize}
\item The channels are NOT in series with the 5 A path (trunk exits the pour at y 43.5;
C1B hangs at y 51, a dead-end stub). Verified geometrically.
\item Worst realistic channel current: PD transition slew (30 mV/us x 10 uF) = 0.3 A for
\textasciitilde{}0.5 ms -\textgreater{} model rise \textasciitilde{}1 C on the 1.11 mm channel alone (3.49 A capacity at 10 C).
Hot-plug edge cases are adiabatic (trivial I2t). No converter on board -\textgreater{} no
sustained ripple duty.
\item Combined 2.36 mm \textgreater{}= the net's own 1.75 mm rule anyway.
\end{itemize}
Acceptable for a bulk cap's duty by \textgreater{}= 10x margin. Recommend recording as a standing waiver so P9's DFM pass does not re-litigate it.

\subsubsection*{Hunt-list verdicts not already covered}
\begin{itemize}
\item 5 A path on render: J1 fan-in solid (pour wraps both merged VBUS pads, solid-connect,
CC-pad ladder slots stay left of x 18.35, \textgreater{} 4 mm continuous corridor); 3.0 mm trunk
continuous (20,43.5)-\textgreater{}(52.2,43.5) straight into J2 pad 1; 0 VBUS vias. PASS.
\item GND under trunk: B.Cu pour solid and unslotted beneath the entire trunk shadow
(verified on raster; one connected component, 1322 mm2). The 9 B.Cu signal crossings
(CC1 x2, CC2 x2, /VIND x1, /VAUX x4) all sit peripheral - CC pair slots the J1 fan
region where the funnel feeds around them from N and S; /VIND slot is 7 mm north of
J2's GND pad; /VAUX slots the far top-left corner. None slot the return. PASS
(E1 is the return's real problem, not the crossings).
\item TVS D1: inside the pour 3.2-3.6 mm from J1's VBUS pads, first element on the net; GND
side = 3 short 0.2/0.3 mm traces into 3 vias into solid pour. Est. clamp-loop L well
under 10 nH; at the design threat's di/dt (\textasciitilde{}5 A/us cable-spike class) that is tens of
mV on top of a 28 V clamp - negligible. Surge I2t through 3 vias: adiabatic, fine.
PASS. Bulk at entry: C1A/C2 tap the trunk 2.5-6.5 mm from the pour, C1B in the pour;
20 uF total \textless{} 100 uF cSnkBulkPd. PASS.
\item User-facing reality: DIP reachable at top edge, actuators clear; J2 wire entry faces
off-board (P6 WRL-verified rotation); J3 pluggable at top-left; LED row visible at
bottom edge. All PASS mechanically - every labeling promise FAILS (E2/E3).
\item Waiver soundness: W1 (DP/DM short at chip) re-verified in copper - U1's DP/DM are
net /BC12\_DIS, J1's A6/A7/B6/B7 pads unconnected; hardware-safer than connector
wiring, no downstream copper effect; waiver stands but its silk condition is
unfulfilled (E2). W2 (VDD worst corner) unaffected by routing; R2A+R2B = 2x 510R
1206 in series confirmed placed 25 mm from U1 (\textgreater{}= 8 mm rule met); bring-up
measurement stands. U1 pad 11 (baseplate, "pin 0" trap V11) confirmed net GND with
3 vias under it.
\item Warnings triage: the machine record contains zero warnings anywhere (ERC 0/0, DRC
0/0 err+warn, verify 8 checks 0 findings). Only tool-emitted advisories: plane\_repair
dead island (A3, triaged) and stitch\_vias relocations (benign, verified placed).
route completion 1.0 at final; the 0.956 probe figure is historical.
\end{itemize}
\subsubsection*{Open (could not judge from artifacts)}
1. SW1 slide direction: whether the physical actuator closes toward the footprint's "ON" label on this exact SHOU HAN part cannot be proven from the board file. Wrong direction inverts the (missing) table's ON column. Fail-safe floor exists (open = 5 V), but add to bring-up: verify one non-5 V profile before trusting the map. 2. Real steady-state temperature of the E1 necks (IPC long-trace model overstates short heatsunk necks; could be 30-60 C rise instead of 133 C). Irrelevant if E1 is fixed - the fix is minutes of work. 3. Order-doc promises (PD-only, e-marked-cable disclaimer, Q1 pinout buzz-out, VDD measurement) - P9/P10 artifacts do not exist yet; carried forward, not judgeable.

\section{DFM and Fabrication}
Gate \texttt{dfm} (phase P9): \textbf{pass} after 2 attempt(s); last 2026-08-14T01:20:22, failing 0 of 20.
DFM findings by severity: warning: 20.
\subsection*{Fabrication outputs}
Exported layers: F.Cu, B.Cu, F.Silkscreen, B.Silkscreen, F.Mask, B.Mask, F.Paste, B.Paste, Edge.Cuts.
\subsection*{Bill of materials}
{\small
\begin{longtable}{p{4.4cm}p{2.8cm}p{4.6cm}l}
\toprule
\textbf{Comment} & \textbf{Designator} & \textbf{Footprint} & \textbf{LCSC} \\
\midrule
\endhead
10uF 50V X5R & C1A,C1B & C1206 & C13585 \\
100nF 50V X7R & C2 & C0603 & C14663 \\
1uF 50V X5R & C5 & C0603 & C15849 \\
TVS2200 22V & D1 & WSON-6\_L2.0-W2.0-P0.65-BL-EP & C523793 \\
BZX84C6V2 & D2 & SOT-23-3\_L2.9-W1.3-P1.90-LS2.4-BR & C841160 \\
LED YG PWR & D3 & LED0603-RD\_GREEN & C965805 \\
LED RED 5V ONLY & D5 & LED0603-RD & C7496813 \\
LED YG PROFILE OK & D6 & LED0603-RD\_GREEN & C965805 \\
PPTC 1A 30V & F1 & F1206 & C5358568 \\
USB-C 16P 5A & J1 & USB-C-SMD\_MC-311D & C5184243 \\
SCREW 5.08 2P & J2 & CONN-TH\_P5.08\_KF128-5.08-2P & C474952 \\
HDR 1x2 AUX & J3 & HDR-TH\_2P-P2.54-V-M-3 & C32713268 \\
BC847BS & Q1 & SOT-363-6\_L2.2-W1.3-P0.65-LS2.1-BL & C41375126 \\
10k & R1,R8 & R0603 & C25804 \\
510R & R2A,R2B & R1206 & C25386 \\
100k & R3,R4,R5 & R0603 & C25803 \\
6.8k & R6 & R0603 & C23212 \\
4.7k & R7 & R0603 & C23162 \\
47k & R9 & R0603 & C25819 \\
3.3k & R10 & R1206 & C137292 \\
1.5k & R12 & R0603 & C22843 \\
4.7k & R13 & R0805 & C17673 \\
0R & R14 & R0603 & C21189 \\
DIP-3 & SW1 & SW-SMD\_6P-L7.6-W6.0-P2.54-LS9.3-BL & C7421520 \\
CH224K & U1 & ESSOP-10\_L4.9-W3.9-P1.0-LS6.0-TL-EP & C970725 \\
\bottomrule
\end{longtable}
}
2 CPL rotation correction(s) applied; no missing LCSC numbers.
\subsection*{Quote}
Estimated total \$19.65 for qty 10 (unit \$1.96). This is an estimate from published pricing; the JLCPCB cart quote page is authoritative.
\subsection*{Human steps before ordering}
\begin{enumerate}
\item API order created: orderId a660e81dac5d4fa58fb58f80741327fd, batchNum W2026073002475378. Poll with order\_track.py --workspace \textless{}ws\textgreater{}; verify payment state in the JLCPCB portal (may auto-deduct JLC Balance).
\item COPPER WAIVER: ordered at 1 oz despite the 2 oz design note (explicit --copper-oz human decision) - current paths sized for 2 oz are derated on this rev
\item BOARD-SPECIFIC before ordering: (1) 2oz copper MUST be selected at JLC order time (stackup JLC2313\_1.6\_2oz - the quote page copper-weight option, NOT default 1oz; the 5A path depends on it). (2) PD SOURCES ONLY + e-marked cable for 5A contracts - user docs. (3) Bring-up: measure /VDD at the 5V profile (CH224K IDD unpublished, W2); buzz Q1 E/B/C once (clone PDF unverifiable); continuity SW1 pins 1-3 unpressed. (4) U1 thermal pad ships 100\% paste (declined windowing - watch for float on first article). (5) Hand-solder J2/J3 THT after SMT.
\item Upload boards\textbackslash{}pd-trigger\textbackslash{}fab\textbackslash{}pd-trigger\_gerbers.zip to https://jlcdfm.com/ and review the DFM report (V6: no public API).
\item Upload the same zip at https://cart.jlcpcb.com/quote and confirm the real price against the estimate above.
\item If assembling: upload BOM.csv and CPL.csv, then CHECK THE RENDERED PART PREVIEW for every polarized part (LED/diode/electrolytic) - rotation corrections are the classic failure mode.
\item Review, then pay. Payment is always the human's action.
\end{enumerate}
\section{Run Record (Appendix)}
\subsection*{Phase digests}
\paragraph*{P0}
\begin{quote}\small\ttfamily\raggedright
\# P0 digest (pd-trigger)\\
Novel brief. \textgreater{}3A safety flag engaged: 6 closed-form questions answered as\\
delegate (5A uniform, PTC, 8A terminal, aux=sense-only, soft size, indicated\\
fallback). 0 human interactions (full-auto directive).
\end{quote}
\paragraph*{P1}
\begin{quote}\small\ttfamily\raggedright
\# P1 digest (pd-trigger)\\
3 researchers (scout sonnet, interface+power opus). Scout: LEAD CH224K\\
(C970725 ESSOP-10, straps verified from datasheet), alt HUSB238A (QFN);\\
IP2721/CYPD3177/FUSB302/CH224Q disqualified with reasons; JLC placeholder-row\\
hazard found. Interface: CC no-impedance (BMC 300k), receptacle 5A =\\
collective 1.25A/contact + 20V-rating trap, TVS MANDATORY (chip-death\\
precedent), VDD shunt-dropper 0.28W trap, DP/DM-short vs check\_diffpair trap\\
(diff\_pairs: {[}{]} required). Power: 2oz REQUIRED (stackups.yaml entry added),\\
main-path PTC dropped (decision), LDO over dropper, VDD IDD unresolved -\textgreater{}\\
extract confirms, thermal entries J1/U2.
\end{quote}
\paragraph*{P2}
\begin{quote}\small\ttfamily\raggedright
\# P2 digest (pd-trigger)\\
6 blocks / 1 flat sheet / \textasciitilde{}32 parts. JLC2313\_1.6\_2oz 2L 45x25. VBUS\\
end-to-end 1.75mm@2oz, no vias in the 5A path. diff\_pairs {[}{]} (short trap),\\
GND structural. DIP 3-pos straps (fail-safe 5V) + 4-LED truth table incl.\\
zener-window "5V ONLY"/"PROFILE OK". Cost \textasciitilde{}\$5-7/board. H1 AUTO-approved.\\
Risk: no series element in 100W path (documented chain D1). P3 bounce-backs:\\
LDO Vin\textgreater{}=30V availability; CH224K IDD placeholder.
\end{quote}
\paragraph*{P3}
\begin{quote}\small\ttfamily\raggedright
\# P3 digest (pd-trigger)\\
parts.json 24 lines / 28 refs (post-amendment); extracts C970725 (CH224K:\\
pin map verified, CFG table verified, VDD shunt + 1k dropper canonical,\\
VBUS sense 13.5V abs max + 10k mandatory, NO CC Rd, IDD unpublished -\textgreater{}\\
bring-up item) + C14170 (L78L33 - later REMOVED by amendment). P2 bounce-\\
back executed: LDO dropped (5 reasons), dropper 2x510R-1206-series 3x\\
derated, straps 100k to /VDD, PG unconnected, R14 0R /VIND stub, LEDs\\
re-specced (green-gap catalog ceiling flagged). Librarian: 24/24 pulled\\
clean first pass, U1 pad-11=GND trap RESOLVED (symbol+footprint agree),\\
J1 land exact vs GCT, 4 approved edits (peg np\_thru x2 + LED dots + SW1\\
internal texts) -\textgreater{} scratch DRC 17-\textgreater{}0; model paths abs-fixed; EDITS.md.\\
SW1 pairs 1-6/2-5/3-4 (generator must hard-code). 14/16 courtyards \textless{} pad\\
field (placelib compensates).
\end{quote}
\paragraph*{P4}
\begin{quote}\small\ttfamily\raggedright
\# P4 digest (pd-trigger)\\
Flat sheet, 30 refs/24 lines, generator + retype (only U1 VDD/GND power\_in).\\
ERC 0/0; audit pass (3 structural power\_no\_consumers on pass-through nets).\\
Agent corrected 3 orchestrator-brief errors from file authorities; edited\\
architecture constraints for R2A/R2B split (retro-approved; finding 31:\\
anneal silently drops absent separation refs). Fable reviewer 0 err/2 warn:\\
DP/DM short-at-chip RULED safer than datasheet wiring (W1, PD-only doc);\\
VDD corner math verified 0.979mA@4.4V (W2, bring-up measure). Window\\
detector verified at 4.4/5.25/6.7/9/20V corners; Q1 pin map from primary\\
source; TVS protects every downstream rating; 5A path has nothing in series.\\
H2 AUTO-approved.
\end{quote}
\paragraph*{P5}
\begin{quote}\small\ttfamily\raggedright
\# P5 digest (pd-trigger)\\
board\_init: 45x25 FAIL (14 silk-edge) -\textgreater{} 48x30 (soft-target +20\%) 11 -\textgreater{} two\\
S14 fixes (kc.py refdes regex now takes R2A/R2B suffixes - cross-part silk\\
was mis-classified single-ref; silk\_edge\_clearance = transient by nature)\\
-\textgreater{} PASS: parity 0, setup 0, 31 transient silk, 68 unconnected. 2 LAYERS at\\
2oz (JLC2313\_1.6\_2oz stackup - S14 reference addition) - capability class\\
2layer\_2oz, 12 rules incl. 4 power widths (VBUS 1.75mm@2oz). Tests +2.
\end{quote}
\paragraph*{P6}
\begin{quote}\small\ttfamily\raggedright
\# P6 digest (pd-trigger)\\
Anneal candidates (hpwl 264) REJECTED - all violated the 5A corridor; agent\\
hand-floorplanned (369.6): J1\textless{}-\textgreater{}J2 facing across a clear 5.5mm-tall corridor,\\
D1 3.6mm from J1, bulk at entry, U1 CC-row toward J1 (5/7mm), droppers 23mm\\
out (8mm floor), SW1 top edge, LED row bottom, window cluster compact.\\
Connector rotations PROVEN by WRL pin-fitting (J2 was seed-backwards).\\
31 transient silks -\textgreater{} 0 via 30 move\_text. Gate PASS 0/0; DRC 68 unconnected\\
only; probe 0.971 - residual = J1 VBUS pad-pair merge needs CC1/CC2 on\\
B.Cu (structural, prescribed to P7). H3 AUTO-approved.
\end{quote}
\paragraph*{P7}
\begin{quote}\small\ttfamily\raggedright
\# P7 digest (pd-trigger)\\
Both prescribed premises FALSIFIED by measurement: 1.75mm tracks cannot\\
touch J1 VBUS pads (1.465 max), CC-displacement wouldn't have helped (pads\\
pinch, not escapes). Solution: F.Cu POUR fan-in/merge (3.5mm copper, zones\\
exempt from track width rules) + 3.0mm trunk + 0 VBUS vias. route\_critical\\
KRT attempt self-rolled-back (9 shorts) - abandoned honestly. rules\_gen\\
netclass defect found (one Power class at max width would force 1.75mm into\\
0.5mm-pitch pads): split in .kicad\_pro, DRU untouched. FR rung 1 completion\\
1.0; dedup removed 3 echoes (S14 fix live); stitch 8 GND vias (2 relocated);\\
plane\_repair 0 splits. C1B pour-channel honesty disclosed (1.11+1.25mm,\\
checker-blind). Gate PASS 0/0. Cleanup skipped per V13.
\end{quote}
\paragraph*{P8}
\begin{quote}\small\ttfamily\raggedright
\# P8 digest (pd-trigger)\\
verify PASS 8/8 first try. Fable reviewer: 4 err + 2 warn. Copper errors\\
REAL: 5A RETURN choked at J1 GND pads (0.2mm necks/0.80A, 1 via) - rebuilt\\
w/ priority-1 F.Cu GND lobes (connect\_pads solid) + 15 new 0.6/0.3 vias -\textgreater{}\\
5.68A/4.70A per pad, 16 vias (D7 \textgreater{}=10 met). Removed 2 stitch vias 0.121mm\\
from THT drills - LIVE-PROVEN KiCad DRC blind spot (via-drill vs same-net\\
pad-drill never checked, finding 34). Silk errors: 17 texts added; agent\\
CORRECTED the orchestrator's B-table spec (raw CFG bits would be inverted\\
vs ON=GND switch reality - printed switch positions). Warns: mouth recess\\
-\textgreater{} order docs; outline 48x30 deviation recorded; C1B channels waived w/\\
numbers. Re-gates: drc\_routed 0/0, verify 8/8. H4 AUTO-approved.
\end{quote}
\paragraph*{P9}
\begin{quote}\small\ttfamily\raggedright
\# P9 digest (pd-trigger)\\
fab\_export 11 files 2L; BOM complete (board LCSC fields), 2 rotation\\
corrections; dfm gate PASS FIRST TRY 0/20 warnings (sliver/silk classes,\\
non-gating). The S14 checker fixes (round caps, sliver band, annular union,\\
board-field BOM) all held on a third, novel board.
\end{quote}
\paragraph*{P10}
\begin{quote}\small\ttfamily\raggedright
\# P10 digest (pd-trigger)\\
Quote qty 10 + order manifest ready\_for\_human. Board-specific human\_steps:\\
2oz MUST be selected at order time (quote assumes standard tiers), PD-only\\
+ e-marked cable docs, VDD/Q1/SW1 bring-up measures, U1 paste watch,\\
J2/J3 hand-solder. H5 AUTO-approved-for-package.\\
RUN (c) COMPLETE: novel brief -\textgreater{} order-ready, all 5 gates green, 0 human\\
interactions mid-run. The novel-part-class path (research -\textgreater{} scout -\textgreater{}\\
extract -\textgreater{} bounce-back -\textgreater{} reconcile) exercised end-to-end with 3 architecture\\
corrections from ground truth.
\end{quote}
\subsection*{Gate attempt history}
{\small
\begin{longtable}{llll}
\toprule
\textbf{Timestamp} & \textbf{Gate} & \textbf{Status} & \textbf{Failing/Total} \\
\midrule
\endhead
2026-07-28T08:49:33 & erc & pass & 0/0 \\
2026-07-28T10:14:30 & place & pass & 0/0 \\
2026-07-28T11:01:44 & drc\_routed & pass & 0/0 \\
2026-07-28T11:04:29 & drc\_routed & pass & 0/0 \\
2026-07-28T11:04:44 & verify & pass & 0/0 \\
2026-07-28T11:53:42 & drc\_routed & pass & 0/0 \\
2026-07-28T11:53:42 & verify & pass & 0/0 \\
2026-07-28T11:55:26 & drc\_routed & pass & 0/0 \\
2026-07-28T11:55:27 & verify & pass & 0/0 \\
2026-07-28T11:56:33 & dfm & pass & 0/20 \\
2026-08-14T01:20:22 & dfm & pass & 0/20 \\
\bottomrule
\end{longtable}
}
\subsection*{Issues}
{\small
\begin{longtable}{llp{5cm}ll}
\toprule
\textbf{Id} & \textbf{Phase} & \textbf{Kinds} & \textbf{Severity} & \textbf{Status} \\
\midrule
\endhead
1 & P8 & current-return & error & fixed \\
2 & P8 & current-return & error & fixed \\
3 & P8 & promised-artifact & error & fixed \\
4 & P8 & promised-artifact & error & fixed \\
5 & P8 & mechanical & warning & fixed \\
6 & P8 & architecture-drift & warning & fixed \\
7 & P8 & current-return & info & fixed \\
8 & P8 & architecture-drift & info & fixed \\
9 & P8 & cleanup & info & fixed \\
10 & P8 & current-margin & info & fixed \\
\bottomrule
\end{longtable}
}
\subsection*{Human checkpoints}
{\small
\begin{longtable}{lllp{2.6cm}p{6.4cm}}
\toprule
\textbf{Id} & \textbf{Phase} & \textbf{Status} & \textbf{When} & \textbf{Note} \\
\midrule
\endhead
1 & P2 & approved & 2026-07-28T06:52:56 & AUTO per directive: 2L 2oz 45x25, \textasciitilde{}\$5-7/board, riskiest = no series protection on the 5A path (charger OCP + TVS + sized copper; chain in decisions.md D1) \\
2 & P4 & approved & 2026-07-28T09:09:11 & AUTO per directive: ERC 0/0, audit pass, fable reviewer 0 err/2 warn waived (DP/DM ruling + VDD corner); PDF reports/schematic.pdf \\
3 & P6 & approved & 2026-07-28T10:14:30 & AUTO per directive; renders reports/pd-trigger\_\{top,bottom,left,right\}.png; J1 flush-mouth + J2 screw headroom flagged for order docs \\
4 & P8 & approved & 2026-07-28T11:55:29 & AUTO per directive: verify 8/8 x2, drc\_routed 0/0; reviewer 4 err ALL FIXED (J1 return 0.8-\textgreater{}5.68/4.70A + 16 vias; 17 silk texts incl. corrected switch-position table), 2 warn (mouth recess -\textgreater{} order docs; outline 48x30 deviation recorded); C1B channels standing waiver. Renders reports/render\_final/ \\
5 & P10 & approved & 2026-07-28T11:56:34 & AUTO per directive: order.json ready\_for\_human; 2oz-selection + PD-only + bring-up steps folded in; payment never automated \\
\bottomrule
\end{longtable}
}
\section{Artifact Index}
{\small
\begin{longtable}{p{9cm}rp{4cm}}
\toprule
\textbf{File} & \textbf{Size} & \textbf{Note} \\
\midrule
\endhead
\texttt{state.json} & 23,483 B &  \\
\texttt{kicad/pd-trigger.kicad\_pcb} & 234,842 B &  \\
\texttt{kicad/pd-trigger.kicad\_sch} & 150,179 B &  \\
\texttt{fab/pd-trigger\_gerbers.zip} & 38,673 B & sha256 ac8e03ac89de8d2b... \\
\texttt{fab/order.json} & 3,717 B &  \\
\texttt{fab/BOM.csv} & 953 B &  \\
\texttt{fab/CPL.csv} & 1,017 B &  \\
\texttt{brief/brief.md} & 832 B &  \\
\texttt{requirements.md} & 11,271 B &  \\
\texttt{architecture/blocks.md} & 13,604 B &  \\
\texttt{architecture/constraints.json} & 3,694 B &  \\
\texttt{architecture/decisions.md} & 37,537 B &  \\
\texttt{architecture/power\_tree.md} & 9,068 B &  \\
\texttt{architecture/sheets.md} & 8,918 B &  \\
\texttt{architecture/stackup.md} & 4,553 B &  \\
\texttt{reports/board\_init.json} & 25,835 B &  \\
\texttt{reports/bom\_cpl.json} & 12,202 B &  \\
\texttt{reports/drc-place.json} & 51,463 B &  \\
\texttt{reports/erc-waivers.md} & 719 B &  \\
\texttt{reports/erc\_p4.json} & 256 B &  \\
\texttt{reports/fab\_export.json} & 2,295 B &  \\
\texttt{reports/fp\_scratch\_drc.json} & 1,121 B &  \\
\texttt{reports/fp\_verify\_U1.json} & 1,575 B &  \\
\texttt{reports/gate-dfm.json} & 1,082 B &  \\
\texttt{reports/gate-drc\_routed.json} & 403 B &  \\
\texttt{reports/gate-erc.json} & 473 B &  \\
\texttt{reports/gate-place.json} & 414 B &  \\
\texttt{reports/gate-verify.json} & 705 B &  \\
\texttt{reports/lib\_pull.json} & 8,948 B &  \\
\texttt{reports/lib\_verify.json} & 22,907 B &  \\
\texttt{reports/netlist\_audit\_p4.json} & 2,097 B &  \\
\texttt{reports/pd-trigger.net} & 65,211 B &  \\
\texttt{reports/pd-trigger\_bottom.png} & 59,102 B &  \\
\texttt{reports/pd-trigger\_left.png} & 27,161 B &  \\
\texttt{reports/pd-trigger\_right.png} & 26,711 B &  \\
\texttt{reports/pd-trigger\_top.png} & 237,347 B &  \\
\texttt{reports/place\_anneal.json} & 2,800 B &  \\
\texttt{reports/place\_edit\_p6.json} & 6,708 B &  \\
\texttt{reports/place\_edit\_silk.json} & 1,968 B &  \\
\texttt{reports/place\_metrics\_final.json} & 5,110 B &  \\
\texttt{reports/place\_p6.ops.json} & 6,476 B &  \\
\texttt{reports/place\_seed.json} & 1,851 B &  \\
\texttt{reports/place\_seed.ops.json} & 3,587 B &  \\
\texttt{reports/plane\_repair\_flag.json} & 1,301 B &  \\
\texttt{reports/planes\_gen.json} & 682 B &  \\
\texttt{reports/planes\_gen\_gnd.json} & 809 B &  \\
\texttt{reports/planes\_gen\_vbus.json} & 596 B &  \\
\texttt{reports/review-board.json} & 6,037 B &  \\
\texttt{reports/review-board.md} & 11,514 B &  \\
\texttt{reports/review-schematic.json} & 2,488 B &  \\
\texttt{reports/review-schematic.md} & 17,762 B &  \\
\texttt{reports/review\_bcu.svg} & 23,093 B &  \\
\texttt{reports/review\_fcu.svg} & 39,795 B &  \\
\texttt{reports/route-probe-deep.json} & 323 B &  \\
\texttt{reports/route-probe.json} & 323 B &  \\
\texttt{reports/route\_auto.json} & 3,405 B &  \\
\texttt{reports/route\_edit\_fix.json} & 1,460 B &  \\
\texttt{reports/route\_edit\_fix2.json} & 256 B &  \\
\texttt{reports/route\_edit\_fix3.json} & 864 B &  \\
\texttt{reports/route\_edit\_gnd\_vias.json} & 1,918 B &  \\
\texttt{reports/route\_edit\_vbus.json} & 955 B &  \\
\texttt{reports/rules\_gen.json} & 1,166 B &  \\
\texttt{reports/schematic.pdf} & 105,717 B &  \\
\texttt{reports/stitch\_dryrun.json} & 3,172 B &  \\
\texttt{reports/stitch\_vias.json} & 3,187 B &  \\
\texttt{reports/top.net} & 65,211 B &  \\
\texttt{reports/U1\_ESSOP10.overlay.svg} & 4,426 B &  \\
\texttt{reports/verify\_all.json} & 1,669 B &  \\
\bottomrule
\end{longtable}
}
\end{document}
